\hypertarget{appendix-280445746493}{}
\section*{Bit significance}
\section{Bit significance}\label{bit-significance}

\textbf{Curriculum depth:} Introduced · \textbf{First seen:} C001 · \textbf{Latest development:} C001

\subsection{Definition and mechanism}\label{definition-and-mechanism}

LSB/MSB describe mathematical weights, not byte order in memory.

LSB means least significant bit, normally position 0 with weight 1. MSB is the highest significant position; in an eight-bit unsigned value it has weight 128. A displayed leftmost bit, first packed array element, and first byte at an address are separate ordering choices.

\subsection{Prerequisites}\label{prerequisites}

\begin{itemize}
\tightlist
\item
  \hyperref[appendix-83c3f28f7bda]{Positional notation} --- A base is the number of available digits.
\end{itemize}

\subsection{Used by}\label{used-by}

\begin{itemize}
\tightlist
\item
  \hyperref[appendix-f414bd4ed05c]{Interpretation contract} --- uses bit significance as part of its representation or reasoning.
\item
  \hyperref[appendix-6b8075763724]{Hexadecimal} --- uses bit significance as part of its representation or reasoning.
\item
  \hyperref[appendix-dbe5905163c2]{Bit mask} --- uses bit significance as part of its representation or reasoning.
\item
  \hyperref[appendix-73425b25118e]{Shift operations} --- uses bit significance as part of its representation or reasoning.
\item
  \hyperref[appendix-bc07c3d25b18]{Endianness} --- uses bit significance as part of its representation or reasoning.
\item
  \hyperref[appendix-ce20506c8532]{Signed magnitude} --- uses bit significance as part of its representation or reasoning.
\end{itemize}

\subsection{Recall and next depth}\label{recall-and-next-depth}

Explain bit significance without the label, work the example above, then state the assumption whose removal would change the result.

\textbf{Pending:} bit numbering conventions and their separation from endianness.

\subsection{Source and navigation}\label{source-and-navigation}

\href{Continuous\%20Notes/Computing\%20Foundations\%20-\%20From\%20Bits\%20to\%20C++\%20and\%20Systems\%20APIs\#From\%20a\%20signal\%20to\%20an\%20interpretation}{C001 teaching chapter} · \hyperref[appendix-fd090d982a15]{Part definitions and mechanisms} · \href{}{Dependency Map (02-Foundations-II, appendix-4ad47bdd269a)}
